% problem-000069 3 化学热力学计算
\section*{3. 化学热力学计算}
\textit{下列为分问。}

\subsection*{(1)}
$\mathrm { F e O ( s ) + C O ( g )  \ F e ( s ) + C O _ { 2 } ( g ) }$ $K _ { 1 } ^ { \ominus } \ = 1 . 0 0 \ \left( T \ = 1 2 0 0 \mathrm { K } \right)$ (2) $\mathrm { F e O ( s ) } + \mathrm { C ( s ) }  \mathrm { F e ( s ) } + \mathrm { C O ( g ) }$ $K _ { 2 } ^ { \ominus }$ (� = 1200 K)

⽓体常数 $R = 8 . 3 1 4 \mathrm { J } \mathrm { m o l } ^ { - 1 } \mathrm { K } ^ { - 1 }$ ；相关的热⼒学数据(298 K)列⼊下表：

\textit{（原卷此处含表格/仪器试剂表，详见 PDF 或 MD 源文件。）}

\subsection*{3-1}
假设上述反应体系在密闭条件下达平衡时总压为1200 kPa，计算各⽓体的分压。

\subsection*{3-2}
计算 $K _ { 2 } ^ { \ominus }$ 。

\subsection*{3-3}
计算 $\mathrm { C O _ { 2 } ( g ) }$ 的标准熵 $S _ { m } ^ { \ominus }$ （单位：J $\mathrm { m o l ^ { - 1 } K ^ { - 1 } } )$ 。设反应的焓变和熵变不随温度变化。

\subsection*{3-4}
反应体系中，若 $\operatorname { C O } ( \mathrm { g } )$ 和 $\mathrm { C O } _ { 2 } ( \mathrm { g ) }$ 均保持标态，判断此条件下反应的⾃发性（填写对应的字⺟）：3-4-1 反应(1)： A. ⾃发 B. 不⾃发 C. 达平衡3-4-2 反应(2)： A. ⾃发 B. 不⾃发 C. 达平衡

\subsection*{3-5}
若升⾼温度，指出反应平衡常数如何变化（填写对应的字⺟）。3-5-1 反应(1)： A. 增⼤ B. 不变化 C. 减⼩3-5-2 反应(2)： A. 增⼤ B. 不变化 C. 减⼩
